\section{Chi phí}

\begin{frame}{Chi phí trong năm đầu}{Các chi phí dự tính ban đầu}

\begin{table}[ht]
\centering
\small
\begin{tabular}{p{4.2cm} r p{6.2cm}}
\hline
\textbf{Hạng mục chi phí} &
\textbf{Mức dự kiến (VNĐ)} &
\textbf{Phân bổ} \\
\hline

\textbf{Marketing} (CAC, Khuyến mãi) &
250.000.000 &
Ads mạng xã hội, voucher giải quyết Cold-start. \\

\textbf{Phát triển Sản phẩm} (Tech/R\&D) &
30.000.000 &
Phát triển tính năng Escrow, tích hợp eKYC. \\

\textbf{Hạ tầng Công nghệ} (Server, API) &
50.000.000 &
Phí Cloud Server, OTP, API đối soát, Domain name,... \\

\textbf{Vận hành \& Nhân sự} (Operations) &
264.000.000 &
Thiết kế, CSKH, bổ sung team Marketing, Finance (Lương mỗi thành viên: 7tr/tháng). \\

\textbf{Tư vấn Pháp lý \& Luật} &
40.000.000 &
Soạn hợp đồng, quy định, bổ sung nhân sự luật. \\

\textbf{Quỹ Dự phòng} &
50.000.000 &
Dự phòng rủi ro tài sản và thanh khoản (khoảng 10\%). \\

\hline
\textbf{TỔNG CHI PHÍ DỰ KIẾN (NĂM 1)} &
\textbf{684.000.000} &
\\
\hline
\end{tabular}
\end{table}

\end{frame}

\section{Doanh thu}

\begin{frame}{Nguồn tiền trên nền tảng năm đầu tiên}
    \framesubtitle{Dự tính tổng giao dịch có được qua từng quý}

     \centering
     \resizebox{.96\textwidth}{!}{%
     \tiny
     \setlength{\tabcolsep}{5pt}
     \renewcommand{\arraystretch}{1.35}

     \begin{tabular}{lrrrrr}
        \hline
        \textbf{Giai đoạn}
        & \textbf{User ví/Tổng User}
        & \textbf{Số đơn}
        & \textbf{GD theo đơn}
        & \textbf{GD mở ví}
        & \textbf{Tổng GD Quý (GMV)}
        \\
        \hline

        Quý 1 (Thử nghiệm)
        & 50/500
        & 150
        & 39.750.000
        & 250.000
        & \textbf{40.000.000}
        \\

        Quý 2 (Tăng trưởng)
        & 100/2.000
        & 600
        & 159.000.000
        & 500.000
        & \textbf{159.500.000}
        \\

        Quý 3 (Chiếm lĩnh)
        & 300/5.000
        & 1.800
        & 477.000.000
        & 1.500.000
        & \textbf{478.500.000}
        \\

        Quý 4 (Mở rộng \& B2B)
        & 600/10.000
        & 4.500
        & 1.192.500.000
        & 3.000.000
        & \textbf{1.195.500.000}
        \\

        \hline

        \textbf{TỔNG NĂM 1}
        & \textbf{600/10.000}
        & \textbf{7.050}
        & \textbf{1.868.250.000}
        & \textbf{5.250.000}
        & \textbf{1.873.500.000}
        \\

         \hline
     \end{tabular}
     }

\end{frame}

\begin{frame}{Doanh thu dự kiến năm đầu tiên}
    \framesubtitle{Doanh thu từ giao dịch và phí mở ví theo từng quý}

     \centering
     \resizebox{.96\textwidth}{!}{%
     \tiny
     \setlength{\tabcolsep}{5pt}
     \renewcommand{\arraystretch}{1.35}

    \begin{tabular}{lrrrr}
        \hline
        \textbf{Giai đoạn}
        & \textbf{GMV}
        & \textbf{Doanh thu đơn}
        & \textbf{Doanh thu mở ví}
        & \textbf{Doanh thu tổng}
        \\
        \hline

        \textbf{Quý 1}
        & 40.000.000
        & 9.450.000
        & 250.000
        & \textbf{9.700.000}
        \\

        \textbf{Quý 2}
        & 159.500.000
        & 37.800.000
        & 500.000
        & \textbf{38.300.000}
        \\

        \textbf{Quý 3}
        & 478.500.000
        & 113.400.000
        & 1.500.000
        & \textbf{114.900.000}
        \\

        \textbf{Quý 4}
        & 1.195.500.000
        & 283.500.000
        & 3.000.000
        & \textbf{286.500.000}
        \\

        \hline

        \textbf{CẢ NĂM}
        & \textbf{1.873.500.000}
        & \textbf{444.150.000}
        & \textbf{5.250.000}
        & \textbf{449.400.000}
        \\

         \hline
     \end{tabular}
     }

\end{frame}

\begin{frame}{500 triệu VNĐ Pre-Seed}{Khoản vốn để vượt qua giai đoạn cold-start}
  \vspace{1.2em}
  \centering
  \metric{500 triệu}{vốn cần gọi (20.000 USD)}\hfill
  \metric{350 triệu}{thiếu hụt lớn nhất ở Q3}\hfill
  \metric{+116 triệu}{dòng tiền hoạt động Q4}\hfill
  \metric{Q3 năm 2}{hòa vốn lũy kế}
  \vfill
  \centering\small\color{cookieMuted}
  350 triệu bù thiếu hụt tiền mặt tối đa; 150 triệu còn lại là bộ đệm vận hành và rủi ro.
\end{frame}

\begin{frame}{Vốn cần bao nhiêu là đủ?}{Con số 500 triệu đến từ đáy dòng tiền tích lũy, không phải ước lượng}
  \vspace{.2em}
  \centering
  \begin{tikzpicture}
    \begin{axis}[
      mutux axis,
      width=.86\textwidth,
      height=.58\textheight,
      ymin=-420, ymax=90,
      enlarge x limits=.14,
      ylabel={Dòng tiền tích lũy (triệu VNĐ)},
      symbolic x coords={Q1,Q2,Q3,Q4},
      xtick=data,
      axis x line=none,
      extra x ticks={Q1,Q2,Q3,Q4},
      nodes near coords,
      nodes near coords style={font=\small\bfseries,text=cookieInk,anchor=north}
      ]
      \addplot[draw=none,fill=mutuxBlue!80] coordinates {(Q1,-161) (Q2,-294) (Q3,-350) (Q4,-234)};
      \draw[mutuxOrange,line width=1pt,dashed]
        (axis cs:Q1,-350) -- (axis cs:Q4,-350)
        node[above left=0mm and 1mm,font=\scriptsize\bfseries,text=mutuxOrange]
        {Đáy tiền mặt $-$350 triệu};
    \end{axis}
  \end{tikzpicture}
  \vfill
  \centering\small\color{cookieMuted}
  Tổng chi năm 1 là 684 triệu, doanh thu 449 triệu; vốn gọi phủ toàn bộ đáy và giữ bộ đệm rủi ro.
\end{frame}

\begin{frame}{Sử dụng vốn để làm gì?}{MVP đã xong nên vốn dồn cho tăng trưởng và vận hành}
  \vspace{.2em}
  \centering
  \begin{tikzpicture}
    \begin{axis}[
      mutux axis,
      xbar,
      width=.74\textwidth,
      height=.58\textheight,
      xmin=0, xmax=215,
      bar width=13pt,
      enlarge y limits=.15,
      xlabel={Triệu VNĐ},
      ymajorgrids=false,
      xmajorgrids=true,
      symbolic y coords={{Dự phòng},{Pháp lý},{Sản phẩm \& hạ tầng},{Vận hành \& nhân sự},{Marketing \& CAC}},
      ytick=data,
      point meta=explicit symbolic,
      nodes near coords,
      nodes near coords style={font=\scriptsize\bfseries,text=cookieInk,anchor=west,xshift=.6mm}
      ]
      \addplot[draw=none,fill=mutuxBlue] coordinates {
        (180,{Marketing \& CAC}) [180 --- 36\%]
        (150,{Vận hành \& nhân sự}) [150 --- 30\%]
        (80,{Sản phẩm \& hạ tầng}) [80 --- 16\%]
        (40,{Pháp lý}) [40 --- 8\%]
        (50,{Dự phòng}) [50 --- 10\%]};
    \end{axis}
  \end{tikzpicture}
  \vfill
  \centering\small\color{cookieMuted}
  Trọng tâm kỹ thuật: eKYC, escrow, PayOS, ví, đối soát và xử lý tranh chấp.
\end{frame}

\begin{frame}{Giải ngân theo cột mốc}{Vốn được mở dần theo tiến độ kiểm chứng}
  \vspace{1em}
  \begin{columns}[T,onlytextwidth]
    \column{.21\textwidth}
      \cookiekicker{Q1 --- 200 triệu}\par\vspace{.4em}
      Công nghệ lõi eKYC--escrow, pháp lý và khởi động cold-start.
    \column{.21\textwidth}
      \cookiekicker{Q2 --- 150 triệu}\par\vspace{.4em}
      Digital ads, trợ giá đơn đầu, CSKH và đo CAC.
    \column{.21\textwidth}
      \cookiekicker{Q3 --- 100 triệu}\par\vspace{.4em}
      Mở rộng trường đại học, gaming cafe và vận hành.
    \column{.21\textwidth}
      \cookiekicker{Q4 --- 50 triệu}\par\vspace{.4em}
      Bổ sung toàn bộ vào quỹ dự phòng rủi ro tài sản.
  \end{columns}
  \vfill
  \centering\small\color{cookieMuted}
  Tỷ lệ giải ngân 40\% -- 30\% -- 20\% -- 10\% theo bốn quý của năm đầu.
\end{frame}

\begin{frame}{Điều kiện chuyển pha}{Ba chỉ số quyết định việc mở khoản giải ngân tiếp theo}
  \vspace{1.2em}
  \begin{columns}[T,onlytextwidth]
    \column{.30\textwidth}
      \cookiekicker{CAC}\par\vspace{.45em}
      Nhỏ hơn hoặc bằng 60.000 VNĐ mỗi người thuê có đơn hoàn tất.
    \column{.30\textwidth}
      \cookiekicker{Tổn thất thiết bị}\par\vspace{.45em}
      Duy trì dưới 1\% trên tổng số đơn thuê.
    \column{.30\textwidth}
      \cookiekicker{Hoàn tất \& tái thuê}\par\vspace{.45em}
      Hai tỷ lệ này phải cải thiện qua từng quý.
  \end{columns}
  \vfill
  \centering\textcolor{mutuxOrange}{\textbf{Không đạt chỉ số thì không mở pha tiếp theo.}}
\end{frame}
